
\textbf{Proposition 4.5. (Comportement fonctoriel).} \emph{Soit \(f:X'\to X\) un morphisme entre deux \(k\)-sch\'emas propres, et soit \(L\) un module inversible sur \(X\). Alors}~:
\begin{enumerate}[label=(\roman*)]
\item \emph{Si \(L\) est \(\tau\)-\'equivalent \`a \(0\) (resp. num\'eriquement \'equivalent \`a \(0\)), alors \(f^*L\) l'est.}
\item \emph{R\'eciproquement, lorsque \(f\) est surjectif, si \(f^*L\) est \(\tau\)-\'equivalent \`a \(0\) (resp. num\'eriquement \'equivalent \`a \(0\)), alors \(L\) l'est.}
\end{enumerate}

En effet, l'assertion (i) (resp. (ii)) pour la \(\tau\)-\'equivalence r\'esulte aussit\^ot de 4.1 (resp. 4.3). D'autre part, soient \(Y'\) une courbe int\`egre ferm\'ee dans \(X'\), et \(Y=f(Y')\) son image, munie de la structure r\'eduite induite. Lorsque \(\dim(Y)=1\), on a la formule de projection (X 4.4.4)
\[
  \tag*{4.5.1.}
  \langle c_1(f^*L)\cdot Y'\rangle
  =[k(Y'):k(Y)]\,\langle c_1(L)\cdot Y\rangle
\]
o\`u \(k(Y')\) (resp. \(k(Y)\)) est le corps de fonctions de \(Y'\) (resp. \(Y\)), et lorsque \(\dim(Y)=0\), on a \((f^*L)_{Y'}\simeq{O}_{Y'}\). Par cons\'equent, si \(L\) est num\'eriquement \'equivalent \`a \(0\), alors \(\langle c_1(f^*L)\cdot Y'\rangle=0\)\,; d'o\`u (i).

R\'eciproquement, lorsque \(f\) est surjective, \'etant donn\'ee une courbe int\`egre ferm\'ee \(Y\) dans \(X\), on peut trouver une courbe int\`egre ferm\'ee \(Y'\) dans \(X'\) telle que \(Y=f(Y')\)~: soient \(y\) le point g\'en\'erique de \(Y\), \(y'\) un point de la fibre \(f^{-1}(y)\) alg\'ebrique sur \(k(y)\), et \(Y'\) l'adh\'erence de \(y'\). Alors 4.5.1 montre que \(\langle c_1(f^*L)\cdot Y'\rangle=0\) implique \(\langle c_1(L)\cdot Y\rangle=0\)\,; d'o\`u (ii).

\textbf{Th\'eor\`eme 4.6. (La caract\'erisation num\'erique de \(\underline{\Pic}^\tau_{X/k}\)).} \emph{Soient \(k\) un corps alg\'ebriquement clos, \(X\) un \(k\)-sch\'ema propre, et \(L\) un faisceau inversible sur \(X\). Alors les conditions suivantes sont \'equivalentes}~:
\begin{enumerate}[label=(\roman*)]
\item \emph{\(L\) est \(\tau\)-\'equivalent \`a \(0\).}
\item \emph{Pour tout faisceau coh\'erent \(F\) sur \(X\), on a}
